% Fig 2.5 — Four properties of Arabic that break lexical retrieval.
%
% New figure, 2026-08-09, for §2.1.5 (chapter2.tex:116-133).
% Spec: research_decisions/CH2_FIGURES_SPEC.md
%
% REVISED 2026-08-09 on Elhaj's review:
%   - Panel 3: the diacritic example was «المَثَانةُ» / «المثانة» (the bladder), taken from
%     chapter4.tex:130. Replaced on his instruction ("this is not a good word") with
%     «الشَّمْس» / «الشمس» (the sun) — neutral, everyday, and it shows three distinct marks
%     (shadda, fatha, sukun) so the contrast is visible at print size.
%     ⚠️ The CAPTION in chapter2.tex was corrected to claim only the ORTHOGRAPHIC example
%     comes from the Chapter 4 failure analysis. The diacritic one no longer does.
%   - Panel 4: was an abstract "dialect vs MSA" statement. Now shows one question rendered
%     in MSA and five regional varieties, per his request to illustrate multiple dialects.
%
% PROVENANCE OF EVERY ARABIC STRING:
%   Panel 1  root and its four derived forms  -> chapter2.tex:122, verbatim
%   Panel 2  إبن الهيثم  vs  ابن الهيثم        -> chapter4.tex:130, verbatim
%                                                (also MIRACL dev query 32, topics.tsv:4)
%   Panel 3  الشَّمْس / الشمس                    -> illustrative (see note above)
%   Panel 4  dialect forms of "what is this?" -> illustrative; the varieties named are the
%            ones listed at chapter2.tex:125
%
% Arabic inside TikZ requires polyglossia + an Arabic font in THIS preamble; the
% \<...> compat macro from 1-main.tex does not exist here, so \textarabic{} is used.
% Arial matches 1-main.tex:17.
\documentclass[border=4pt]{standalone}
\usepackage{fontspec}
\IfFontExistsTF{Times New Roman}{\setmainfont{Times New Roman}}{\setmainfont{TeX Gyre Termes}}
\usepackage{polyglossia}
\setdefaultlanguage{english}
\setotherlanguage{arabic}
\newfontfamily\arabicfont[Script=Arabic,Scale=1.2]{Arial}
\usepackage{tikz}
% Shared TikZ style for thesis figures.
% Each fig_*.tex \input{} this file after \usepackage{tikz}.
% Defines a muted, examiner-friendly color palette and semantic box styles.

\usepackage{xcolor}
\usepackage{fontawesome5}
\usetikzlibrary{positioning, arrows.meta, shapes.geometric, calc, fit, backgrounds}

% ─── Palette ───────────────────────────────────────────────────────────
% Each role has a stroke color (dark) and a fill color (very light).
% Designed to be grayscale-safe — lightness differs enough that B&W printing
% still preserves the distinction.

\definecolor{thQueryS}{HTML}{4A6FA5}    \definecolor{thQueryF}{HTML}{E8EEF6}   % query / input
\definecolor{thLLMS}{HTML}{6D4A8F}      \definecolor{thLLMF}{HTML}{ECE5F2}     % LLM / generation
\definecolor{thRetS}{HTML}{3D7A4F}      \definecolor{thRetF}{HTML}{E4EFE7}     % retriever
\definecolor{thDataS}{HTML}{6A6A6A}     \definecolor{thDataF}{HTML}{F1F1F1}    % corpus / data
\definecolor{thFusionS}{HTML}{A05A30}   \definecolor{thFusionF}{HTML}{F4E8DF}  % fusion
\definecolor{thOutS}{HTML}{1F6F8A}      \definecolor{thOutF}{HTML}{DEEAEF}     % output / final
\definecolor{thHiS}{HTML}{0D4D63}       \definecolor{thHiF}{HTML}{C8D9DF}      % highlighted (our contribution)
\definecolor{thMuted}{HTML}{5C5C5C}

% ─── Box styles ────────────────────────────────────────────────────────
% Usage in tikzpicture options:  query/.style={thQuery}, llm/.style={thLLM}, ...
% Or directly on nodes:          \node[thLLM] (m) {...};
\tikzset{
  thBase/.style={rectangle, draw, rounded corners=2pt,
                 minimum width=28mm, minimum height=10mm,
                 align=center, font=\small, line width=0.6pt,
                 inner sep=4pt},
  thQuery/.style={thBase, draw=thQueryS, fill=thQueryF},
  thLLM/.style={thBase, draw=thLLMS, fill=thLLMF},
  thRet/.style={thBase, draw=thRetS, fill=thRetF},
  thData/.style={cylinder, shape border rotate=90, aspect=0.3, draw=thDataS, fill=thDataF,
                 minimum width=24mm, minimum height=12mm, align=center, font=\small, line width=0.6pt},
  thFusion/.style={thBase, draw=thFusionS, fill=thFusionF},
  thOut/.style={thBase, draw=thOutS, fill=thOutF},
  thHi/.style={thBase, draw=thHiS, fill=thHiF, line width=1pt},
  thArrow/.style={-{Stealth[length=2.5mm]}, line width=0.7pt, draw=thMuted},
  thLabel/.style={font=\footnotesize\itshape, color=thMuted},
  thStage/.style={font=\bfseries\footnotesize, color=thMuted},
}

% ─── Icon helpers ──────────────────────────────────────────────────────
% Tiny FontAwesome icons that can sit inside a node prefix.
% Usage:  \node[thQuery] {\thIconQuery~User query};
\newcommand{\thIconQuery}{\textcolor{thQueryS}{\faSearch}}
\newcommand{\thIconUser}{\textcolor{thQueryS}{\faUser}}
\newcommand{\thIconLLM}{\textcolor{thLLMS}{\faBrain}}
\newcommand{\thIconRet}{\textcolor{thRetS}{\faNetworkWired}}
\newcommand{\thIconCorpus}{\textcolor{thDataS}{\faDatabase}}
\newcommand{\thIconFusion}{\textcolor{thFusionS}{\faCodeBranch}}
\newcommand{\thIconOut}{\textcolor{thOutS}{\faFile}}
\newcommand{\thIconEval}{\textcolor{thOutS}{\faChartBar}}
\newcommand{\thIconStar}{\textcolor{thHiS}{\faStar}}
\newcommand{\thIconCog}{\textcolor{thMuted}{\faCog}}


\newcommand{\ar}[1]{\textarabic{#1}}

\begin{document}
\begin{tikzpicture}[
    panel/.style={rectangle, draw=thDataS, fill=white, rounded corners=2pt,
                  line width=0.6pt, align=center, inner sep=5pt,
                  minimum width=68mm, anchor=north, font=\small},
    ptitle/.style={font=\bfseries\footnotesize, color=thQueryS, align=center, anchor=south},
    consequence/.style={rectangle, draw=thFusionS, fill=thFusionF, rounded corners=2pt,
                        line width=0.8pt, align=center, inner sep=6pt, anchor=north,
                        minimum width=148mm, font=\small},
    remedy/.style={thBase, draw=thHiS, fill=thHiF, line width=1pt, anchor=north,
                   minimum width=148mm, font=\small}
]

  % ─── Row 1 ───────────────────────────────────────────────────────────
  % Both row-1 panels share a minimum height so the row bottom is deterministic
  % and row 2 can be placed at a fixed y regardless of content length.
  \node[panel, minimum height=30mm] (p1) at (-38mm, 0mm) {%
    \ar{ك\,ت\,ب} \;\textit{(root)}\\[3pt]
    \begin{tabular}{@{}r@{\;\;}l@{\qquad}r@{\;\;}l@{}}
      \ar{كِتاب} & \textit{book}   & \ar{مكتبة} & \textit{library}\\
      \ar{كاتب}  & \textit{writer} & \ar{مكتوب} & \textit{written}\\
    \end{tabular}\\[3pt]
    {\footnotesize one root, dozens of surface forms}};
  \node[ptitle] at (p1.north) [yshift=1.5mm] {1. Morphological richness};

  \node[panel, minimum height=30mm] (p2) at (38mm, 0mm) {%
    \ar{إبن الهيثم}\\[2pt]
    \textit{vs}\\[2pt]
    \ar{ابن الهيثم}\\[3pt]
    {\footnotesize same person, two spellings of Alif}};
  \node[ptitle] at (p2.north) [yshift=1.5mm] {2. Orthographic variation};

  % ─── Row 2 ───────────────────────────────────────────────────────────
  % Reframed 2026-08-09 on Elhaj's instruction: "make it a word that has multiple
  % meanings without diacritics". This is the INVERSE failure of panels 1, 2 and 4 —
  % one surface form standing for several concepts, rather than one concept spread
  % over several forms — and it is what chapter2.tex:131 actually describes
  % ("identical consonantal skeletons may correspond to multiple words with different
  % meanings"). The consequence bar below was widened to state both directions.
  \node[panel, minimum height=50mm] (p3) at (-38mm, -40mm) {%
    \ar{علم} \quad\textit{as written in the corpus}\\[4pt]
    \begin{tabular}{@{}r@{\;\;}l@{}}
      \ar{عِلْم}  & \textit{knowledge}\\[1pt]
      \ar{عَلَم}  & \textit{flag}\\[1pt]
      \ar{عَلِمَ} & \textit{he knew}\\
    \end{tabular}\\[4pt]
    {\footnotesize one spelling, three different words;}\\
    {\footnotesize the diacritics that separate them are}\\
    {\footnotesize omitted in almost all modern text}};
  \node[ptitle] at (p3.north) [yshift=1.5mm] {3. Diacritic sensitivity};

  \node[panel, minimum height=50mm] (p4) at (38mm, -40mm) {%
    \textit{one question, six ways to ask it}\\[3pt]
    \begin{tabular}{@{}l@{\;\;}r@{}}
      \footnotesize MSA (the corpus) & \ar{ما هذا؟}\\[1pt]
      \footnotesize Sudanese         & \ar{شنو ده؟}\\[1pt]
      \footnotesize Egyptian         & \ar{إيه ده؟}\\[1pt]
      \footnotesize Levantine        & \ar{شو هاد؟}\\[1pt]
      \footnotesize Gulf             & \ar{شنو هذا؟}\\[1pt]
      \footnotesize Maghrebi         & \ar{آش هادا؟}\\
    \end{tabular}\\[3pt]
    % NOT "no dialect form shares a term with MSA" — the Gulf form شنو هذا؟ does
    % share هذا. The accurate and sharper claim is about the interrogative, which
    % differs in every single variety and is the discriminative term for retrieval.
    {\footnotesize the interrogative differs in every variety;}\\
    {\footnotesize none of them matches the MSA \ar{ما}}};
  \node[ptitle] at (p4.north) [yshift=1.5mm] {4. Diglossia};

  % ─── Consequence ─────────────────────────────────────────────────────
  % Panel 3 now runs in the opposite direction to 1, 2 and 4, so the consequence must
  % state both failures. Claiming only "one concept -> many forms" would no longer
  % cover panel 3, and claiming it did would be false.
  \node[consequence] (cons) at (0mm, -96mm) {%
    \textbf{surface form and meaning do not correspond one to one}\\[3pt]
    {\footnotesize many forms for one concept \;(1,\,2,\,4)\;
      $\rightarrow$\; the relevant passage is \textbf{missed}}\\[1pt]
    {\footnotesize one form for many concepts \;(3)\;
      $\rightarrow$\; irrelevant passages are \textbf{matched}}};

  \draw[thArrow] (p1.west) -- ++(-5mm,0) |- (cons.west);
  \draw[thArrow] (p2.east) -- ++(5mm,0)  |- (cons.east);
  \draw[thArrow] (p3.south) -- ++(0,-3mm) -| (cons.north);
  \draw[thArrow] (p4.south) -- ++(0,-3mm) -| (cons.north);

  % ─── Remedy ──────────────────────────────────────────────────────────
  \node[remedy] (rem) at ($(cons.south)+(0,-8mm)$) {%
    \textcolor{thHiS}{\faStar}~~remedy applied in this thesis: enrich the query
    \emph{before} retrieval (Section 2.1.4)};
  \draw[thArrow] (cons.south) -- (rem.north);

\end{tikzpicture}
\end{document}
